% status: FINAL 2024-2025
% v20240829

\documentclass[12pt,a4paper,english]{article}

%--- Styling preamble ---

%% Essential

\usepackage[T1]{fontenc} % essential for hyphenation and copy/pasting in the final pdf
\usepackage{babel} % english-language support and hy-phen-a-tion. supports many languages. best practice is to set the language as an option to the \documentclass
\usepackage{graphicx} % enhanced support for graphics
\usepackage{amsmath} % essential additional math options
\usepackage{xcolor} % for more color options, and colored text
\usepackage{subcaption}
\usepackage{booktabs}
\usepackage{longtable}

\usepackage[style=ieee]{biblatex}
\addbibresource{references.bib}

%% Recommended packages

\usepackage[margin=2.5cm]{geometry} % to set page margins
\usepackage{comment} % To enable comment blocks that are not interpreted
\usepackage{newtxtext,newtxmath} % changes fonts for better on-screen reading


%% Essential in Physics Lab: Skills

\usepackage{siunitx} % for appropriate value + unit + error formatting
    \sisetup{separate-uncertainty} % write error out with ±
\usepackage{titling} % for title page with abstract
\usepackage[nottoc]{tocbibind} % bibliography in ToC

\usepackage[hyperindex,colorlinks,linktoc=all]{hyperref} % for clickable links in pdf. essential feature for the TAs! should often be the last package (with the exceptions below)
\usepackage{caption} % to point hyperref links to top of figure/table. put after hyperref
\usepackage{bookmark} % to improve hyperref pdf bookmarks. put after hyperref
\usepackage[nameinlink]{cleveref} % Better linking, e.g. to multiple figures and tables. Goes after hyperref
    \urlstyle{same}
    \hypersetup{ % hyperref/bookmark/cleveref package setup parameters
            colorlinks=true,
            linkcolor=purple,
            filecolor=teal,
            urlcolor=olive,
            citecolor=teal,
            bookmarksopen=true,
            bookmarksopenlevel=0,
            pdflang=en,
            pdfpagemode=UseOutlines,
            pdfstartview=FitB,
            pdfpagelayout=OneColumn,
            pdftitle={Experiment X Report},
            pdfauthor={T Boermsa, J.W. Sprietsma, M.H.O Persson},
        }

% Stuff for code & algorithms
\usepackage{listings}

\usepackage[utf8]{inputenc}
\usepackage[T1]{fontenc}
\usepackage{listings}
\usepackage{xcolor}
 
\definecolor{codegreen}{rgb}{0,0.6,0}
\definecolor{codegray}{rgb}{0.5,0.5,0.5}
\definecolor{codepurple}{rgb}{0.58,0,0.82}
\definecolor{backcolour}{rgb}{0.95,0.95,0.92}
 
\lstdefinestyle{python}{
    backgroundcolor=\color{backcolour},   
    commentstyle=\color{codegreen},
    keywordstyle=\color{magenta},
    numberstyle=\tiny\color{codegray},
    stringstyle=\color{codepurple},
    basicstyle=\footnotesize,
    breakatwhitespace=false,         
    breaklines=true,                 
    captionpos=b,                    
    keepspaces=true,                 
    numbers=left,                    
    numbersep=5pt,                  
    showspaces=false,                
    showstringspaces=false,
    showtabs=false,                  
    tabsize=2
}
\lstset{style=python}

\usepackage{algorithm}

%% Packages not in this example, but that you might want to use:


% General

\usepackage{enumitem} % improve enumeration lists
\usepackage{isodate} % for date formatting
\usepackage{nameref} % to refer to section names themselves, not the number
\usepackage{multirow}


% Figure options

\usepackage{subcaption} % sub-captions. put after package "caption"
\usepackage{wrapfig}
\usepackage{float}


% Math

\usepackage{bm} % for bold face in math environment (recommended by AMS style guide)
\usepackage{mathtools} % for definition symbol :=

% Options for Hyperref and related packages that were used in 2024 in the manuals

\urlstyle{same}
\hypersetup{ % bookmark/hyperref package setup parameters
        colorlinks=true,
        linkcolor=purple,
        filecolor=teal,
        urlcolor=olive,
        citecolor=teal,
        bookmarksopen=true,
        bookmarksopenlevel=0,
        pdflang=en,
        pdfpagemode=UseOutlines,
        pdfstartview=FitB,
        pdfpagelayout=OneColumn,
    }
\hypersetup{ % PDF metadata
    pdftitle={NLP Assignment2},
    pdfauthor={Teun Boersma, Julian Sprietsma and Marcus Persson},
    }


%--- End of styling preamble ---


% General info for title page
\title{
    \vspace*{-2cm}
    
\includegraphics[width=\textwidth]{./img/rugr_fse_logoen_rood_cmyk.eps} \\[1.5cm]
   \\ % insert title here
    {\Large Natural Language Processing \\ }
    {\large{Assignment 2}}
}
\author{Teun Boersma (S5195179),\\ Marcus Persson (S5343798),\\ Julian Sprietsma (S5096219)}
\date{\today}


% Let's go!
\begin{document}


% Title page
\begin{titlingpage}   
    \maketitle
\end{titlingpage}


% Table of Contents (ToC)
{   
    \thispagestyle{empty} % no page number for ToC
    \hypersetup{linkcolor=black} % black links for ToC aesthetics
    \tableofcontents
    \clearpage
}

\section{Dataset \& Setup}
This project focuses on the AGNews Dataset, created by Zhang et al. \cite{AGNews}. The dataset consists of 120k training datapoints and 7.6k testing points. Each datapoint consists of a title and a description of a news article. These are around 45 and 200 characters long on average. Each article is categorised based on topic. The set of labels consists of World, Sports, Business and Sci/Tech ($C=4$, $y\in\{1,2,3,4\}$). This project aims to investigate the performance of baseline models on this task. For hyperparameter tuning and testing, a random development split of 10\% was taken from the training set.

To train models, first, the data needs to be preprocessed. The first step was to concatenate the title and the description into a single string of text, which makes it easier to use both features. After this, some basic preprocessing is done, removing short words 

Next, an embedding needs to be made. For this purpose, a Word2Vec vectorizer was trained on the training data, after filtering out words that appear less than 3 times. An embedding dimension of 100 was used; the vectorizer was trained using a window of length 5 across 10 epochs. 

Every process is carried out with a set random seed; 33 was chosen for this.

\section{Models}
This report details an ablation study on sequence length. This is done for two types of models: a recurrent Long Short-Term Memory (LSTM) Model and a Convolutional Neural Network. Both models are used to encode the sequence, after which a linear model extracts the prediction. Detailed in Equation \eqref{eq:linmodel}. The exact dimensions of $\boldsymbol{W}$ differ for the LSTM and CNN models.

\begin{equation}
    \boldsymbol{a} = \boldsymbol{W}\boldsymbol{x}
    \label{eq:linmodel}
\end{equation}

This results in a vector of 4 logits. To convert these logits to a probability, the softmax activation function is used, specified in Equation \eqref{eq:softmax} \cite{pml1Book}.

\begin{equation}
    p(y=c|\boldsymbol{x}, \boldsymbol{W}) = \frac{e^{a_c}}{\sum_{c'=1}^C e^{a_{c'}} }
    \label{eq:softmax}
\end{equation}

\subsection{CNN}
The architecture used for the CNN focuses on a 1D convolution operation across the sequence. This is done for five filters of different lengths, from 3-7 words long. The convolution operation results in a new sequence of vectors $\boldsymbol{h}(t)$ where $t$ specifies the index in the sequence \cite{Goodfellow-et-al-2016}.

\begin{equation}
    \boldsymbol{h}(t) = (\boldsymbol{x} * \boldsymbol{w})(t) = \sum_n x(t+n)w(t)
\end{equation}

Thus, after applying each of the five convolutions, we are left with five sequences, which are pooled together using a max-pooling operation to five vectors of dimension 100; the ReLU activation function (Equation \eqref{eq:relu}) is called on each. Lastly, these vectors are concatenated and form the input for the final linear model in Equation \eqref{eq:linmodel}.

\begin{equation}
    \text{ReLU} (x) = \max(0,x)
    \label{eq:relu}
\end{equation}

The weights of the CNN were initialised using Xavier Uniform initialisation. Here, weights are sampled from a uniform distribution with bounds based on the model complexity. The exact architecture is presented in detail in Table \ref{tab:cnn-architecture} of Appendix \ref{ap:arch}.

\subsection{LSTM}
The LSTM is a gated Recurrent Neural Network (RNN). The LSTM, aside from the standard hidden state vector $\boldsymbol{h}_t$, also has a cell state $\boldsymbol{c}_t$. Like in many RNNs the input $\boldsymbol{x}_t$ is combined with the last cell state $\boldsymbol{h}_{t-1}$ to provide the input for that timestep $t$. An LSTM sets itself apart by its usage of gates, which determine what information from the cell state is forgotten, what information is added and what information is used to construct the hidden state. These gates are labeled $\boldsymbol{f}_t$, $\boldsymbol{i}_t$ and $\boldsymbol{o}_t$, respectively. This results in a set of update equations shown below \cite{jm3}:

\begin{align}
    \boldsymbol{f}_t &= \sigma(\boldsymbol{U}_f\boldsymbol{h}_{t-1} + \boldsymbol{W}_f\boldsymbol{x}_t) \\
    \boldsymbol{i}_t &= \sigma(\boldsymbol{U}_i\boldsymbol{h}_{t-1} + \boldsymbol{W}_i\boldsymbol{x}_t) \\
    \boldsymbol{o}_t &= \sigma(\boldsymbol{U}_o\boldsymbol{h}_{t-1} + \boldsymbol{W}_o\boldsymbol{x}_t) \\
    \boldsymbol{g}_t &= \tanh(\boldsymbol{U}_g\boldsymbol{h}_{t-1} + \boldsymbol{W}_g\boldsymbol{x}_t) \\
    \boldsymbol{c}_t &= \boldsymbol{c}_{t-1} \odot \boldsymbol{f}_t + \boldsymbol{g}_t\odot \boldsymbol{i}_t \\
    \boldsymbol{h}_t &= \boldsymbol{o}_t \odot \tanh(\boldsymbol{c}_t)
\end{align}

Where $\boldsymbol{U}$ and $\boldsymbol{W}$ are matrices and $\odot$ signifies the elementwise (Hadamard) product. The sigmoid function $\sigma(x)$ is given by Equation \eqref{eq:sigmoid}.

\begin{equation}
    \sigma(x) = \dfrac{e^x}{1+e^x}
    \label{eq:sigmoid}
\end{equation}

In this report, a bidirectional version of this system is used, which entails that the input signal is fed through the system both forwards and backwards, and the hidden states are concatenated. Additionally, a layered structure is used. This means that four LSTM-Cells as detailed above were used. The hidden cells $h_{t}^i$ of the $i$-th layer served as the input ($x_t$ in the equations above) for the  $i+1$'th layer. This structure is capable of capturing complex patterns more easily. The model used in the ablation study consisted of four hidden layers. Due to the increased complexity of the LSTM cells, this was chosen to provide a similar model complexity as the CNN. A more detailed tabular overview of the model architecture is presented in Table \ref{tab:lstm-architecture} of Appendix \ref{ap:arch}

\section{Training Protocol}
The models discussed above were implemented using PyTorch \cite{paszke2019pytorchimperativestylehighperformance}. As briefly mentioned before, the purpose of this report is to investigate the performance of the LSTM and CNN models for a varying sequence length. To do this, a model will be trained on data for each maximum sequence length $L \in \{32,64,128,256\}$ as part of the ablation study; with regard to results reported in the following section, we train with a maximum sequence length $L=256$. The sequences will be padded or cut down to the required length for ease of computation. Both models are trained using the same Word2Vec embedding. The model architectures discussed in Appendix \ref{ap:arch}, will remain constant. Additionally, for both models and all values $L$, the Adam optimiser is used with a starting learning rate of $10^{-3}$ and the cross-entropy loss function is applied. Lastly, training is done for a maximum of 20 epochs, with early stopping configured to stop training if the validation loss does not improve for 3 consecutive epochs. The model with the best validation loss is restored at the end of training.

\section{Results}

\subsection{CNN}

The CNN model reported, for the development set, an accuracy of $0.8869$ and a weighted F1-score of $0.8867$; for the test set, the model reported an accuracy of $0.9178$ and a weighted F1-score of $0.9177$. The model achieved a slightly higher performance on the test set than the development set. The confusion matrix comparison is shown in Table \ref{tab:cnn_cm}. This is discussed in detail in the Error Analysis section. 

\begin{table}[h]
    \centering
    \begin{tabular}{|l|cccc|cccc|}
    \hline
          & \multicolumn{4}{c|}{\textbf{Dev Set}} & \multicolumn{4}{c|}{\textbf{Test Set}} \\ \hline
        \textbf{Pred \textbackslash{} Act} & \textbf{World} & \textbf{Sports} & \textbf{Business} & \textbf{Sci/Tech} & \textbf{World} & \textbf{Sports} & \textbf{Business} & \textbf{Sci/Tech} \\ \hline
        \textbf{World}    & 2599 & 148  & 154  & 75   & 1734 & 36   & 82   & 48   \\
        \textbf{Sports}   & 61   & 2708 & 10   & 10   & 18   & 1867 & 10   & 5    \\
        \textbf{Business} & 102  & 18   & 2652 & 267  & 41   & 12   & 1700 & 147  \\
        \textbf{Sci/Tech} & 67   & 34   & 263  & 2832 & 39   & 9    & 109  & 1743 \\ \hline
    \end{tabular}
    \caption{CNN Confusion Matrix Comparison}
    \label{tab:cnn_cm}
\end{table}

With early stopping, the model finishes training for an optimal state after 5 epochs, resulting in a training loss of $0.1845$ and evaluation loss of $0.2615$. The complete training history is shown in Figure \ref{fig:cnn_loss}.

\begin{figure}
    \centering
    \includegraphics[width=0.8\linewidth]{reports/assignment2/img/cnn_training_history.png}
    \caption{CNN Training History Comparing Training and Evaluation Loss, the evaluation loss is only computed after every full epoch.}
    \label{fig:cnn_loss}
\end{figure}

\subsection{LSTM}

The LSTM model reported, for the development set, an accuracy of $0.8992$ and a weighted F1-score of $0.8990$; for the test set, the model reported an accuracy of $0.9268$ and a weighted F1-score of $0.9268$. The model achieved a slightly higher performance on the test set than the development set. The confusion matrix comparison is shown in Table \ref{tab:lstm_cm}. This is again discussed in detail in the Error Analysis section. 

\begin{table}[h]
    \centering
    \begin{tabular}{|l|cccc|cccc|}
    \hline
          & \multicolumn{4}{c|}{\textbf{Dev Set}} & \multicolumn{4}{c|}{\textbf{Test Set}} \\ \hline
        \textbf{Pred \textbackslash{} Act} & \textbf{World} & \textbf{Sports} & \textbf{Business} & \textbf{Sci/Tech} & \textbf{World} & \textbf{Sports} & \textbf{Business} & \textbf{Sci/Tech} \\ \hline
        \textbf{World}    & 2617 & 166  & 120  & 73   & 1741 & 49   & 68   & 42   \\
        \textbf{Sports}   & 84   & 2676 & 14   & 15   & 30   & 1844 & 19   & 7    \\
        \textbf{Business} & 171  & 15   & 2607 & 246  & 67   & 6    & 1682 & 145  \\
        \textbf{Sci/Tech} & 86   & 33   & 334  & 2743 & 52   & 10   & 130  & 1708 \\ \hline
    \end{tabular}
    \caption{LSTM Confusion Matrix Comparison}
    \label{tab:lstm_cm}
\end{table}

With early stopping, the model finishing training for an optimal state after 11 epochs, resulting in a training loss of $0.2253$ and evaluation loss of $0.2196$. The complete training history is shown in Figure \ref{fig:lstm_loss}.

\begin{figure}
    \centering
    \includegraphics[width=0.8\linewidth]{reports/assignment2/img/lstm_training_history.png}
    \caption{LSTM Training History Comparing Training and Evaluation Loss}
    \label{fig:lstm_loss}
\end{figure}

\section{Error Analysis}
% Same as before + comparison. 
% Update error analysis: ≥10 errors, highlight differences vs Assignment 1 baseline failures.

\subsection{CNN}
Looking at the confusion matrix for the Test set of the CNN, it becomes clear that the most commonly made mistake is that the Business and Sci/Tech classes are confused with each other. This classification error also occurred for the baseline models in Assignment 1. The baseline models also showed a significant amount of classification errors between the categories World and Sports, the more complex CNN, however, shows a lower amount of errors between these categories and instead shows a higher error rate between the Business and World classes. The fact that the World - Sports classification error occurs less when it comes to the CNN could be explained by the fact that the context around a word matters more for the CNN. In assignment 1, it was made clear that specific sport-related terms could lead to a classification in the Sports class, while that may not be correct, including context for these examples leads to a lower error rate. \ref{tab:misclassificationCNN} mentions the 10 lowest confidence errors, taking a look at each of them shows us the following. 

In references 1 and 2, the actual class is Sports, but the predicted class is Sci/Tech. Since the CNN uses filters from 3 - 7 words long, we need to look at combinations of words that can cause this misclassification. For reference 1, this could be "successful US debut by beating Literacy". For reference 2, it could be "standing virtually in the middle". Reference 3 mentions "Just begin Glasgow, Cambridge and Cardiff", which is a way of wording that might occur in sports articles, causing a wrong Sports classification. Reference 4 mentions "Cruise line" and "service problems", which could be confused as Sci/Tech sentences rather than business-related topics. For references 5 and 9, there is mostly a mention of a specific person or a party/event. These examples do not really contain any specific sports terms; they are being wrongly classified as the World class. Reference 6 uses the words "buzz" and "cacophony", which may be the reason for a misclassification as Sci/Tech instead of Business. Sentences like "Planning to proceed with" and "gets official go-ahead of the votes" in reference 7 and 8 very likely appear more in business-related articles than sports-related articles, causing this classification error. Reference 10 is about a water crisis and does not have any mention of Sci/Tech-related terms, so a misclassification as World is not surprising. 


\subsection{LSTM}
For the LSTM a similar effect can be seen. Sci/Tech and Business are the most often confused. In comparison to the baseline models, the amount of misclassifications between Sports and World has also decreased for the LSTM. The second most made classification error is also between the classes Business and World. Just like for the CNN it can be stated that an LSTM relies on context and previous inputs. This means that individual words matter less in the classification and instead word combinations and specific ways of wording sentences matters more. We take a look at \ref{tab:misclassificationLSTM} for 10 specific cases of misclassification and why they could be occurring. 

References 1 and 5 misclassify a Business article as a Sports article. For reference 1, this could be caused by the sentence "downgrading of San Diego's credit rating" while for reference 5, the cause may be a sentence like "The green is still there" or "Stars and Stripes now waves in the background". For reference 2, it is quite difficult to see why the misclassification occurs. The article talks about a lawsuit regarding a governor violating the separation of powers, which does not necessarily sound like a business-related article. We were unable to find a plausible explanation for this missclassification. 

References 3 and 4 have the predicted class Sci/Tech, while the actual class is World. Reference 3 features sentences such as "a guide to the primary get ready", while reference 4 mentions a "Potemkin village Gongzhong". Reference 6 is about a "plan to build a multi-million-dollar race track" and a bid that has advanced far; these are sentences likely to occur in Business articles and less likely in Sport articles, thus causing misclassification. References 7 and 10 misclassify a Sci/Tech article as a Business article, thus we are once again looking for sentences that would regularly occur in business-related articles. For reference 7, this could be "Fabled Hollywood producer and man about town".

For reference 10 it is likely "TravelWatch column" or "Aborigine-owned tour company". Reference 8 does not necessarily contain any sentences that would cause a classification as Sci/Tech. There are no mentions of any technological or scientific phenomena and instead only mentions of lacking funds and world population growth, which is more likely to fall into the World class. Lastly, reference 9 does not contain any Business related words or sentences and instead feautures mentions of "The U.S. Coast Guard" and "States of Alabama, Florida and Mississipi", this causes the LSTM to misclassify the article under the World class.

\subsection{Comparison to Assignment 1}
With respect to Assignment 1, it becomes clear that the misclassifications occur less on singular words, but rather on the context that the sentences appear in. This is logical considering that both the CNN and the LSTM rely on the context of the sentence. The CNN does this using a window over multiple words, while the LSTM remembers or forgets previous words and uses them as context. It can also be seen that the accuracies are higher for the more complex models than the baseline models, as is to be expected. 

\section{Ablation Study Concerning Max Sequence Length}

An ablation study was run concerning the max sequence length with the following values: 16, 32, 64, 128, and 256 tokens. With the configuration, we followed the training protocol as standard, with only the max sequence length being changed with the dataset. The initial hypothesis here is that longer sequences provide the model with more context and we expect a higher accuracy. Table \ref{tab:seq_length_comparison} displays the accuracy results per model per dataset per max sequence length, and Figures \ref{fig:cnn_maxlen} and \ref{fig:lstm_maxlen} display the training history per model across all tested max sequence lengths. 

\begin{table}[h]
    \centering
    \begin{tabular}{|c|c|c|c|}
        \hline
        \textbf{Max Sequence Length} & \textbf{Model} & \textbf{Accuracy (Dev Set)} & \textbf{Accuracy (Test Set)} \\ \hline
        \multirow{2}{*}{16}  & CNN  & 0.8721 & 0.8977 \\ \cline{2-4} 
                             & LSTM & 0.9038 & 0.9300 \\ \hline
        \multirow{2}{*}{32}  & CNN  & 0.8858 & 0.9133 \\ \cline{2-4} 
                             & LSTM & 0.9011 & 0.9255 \\ \hline
        \multirow{2}{*}{64}  & CNN  & 0.8888 & 0.9113 \\ \cline{2-4} 
                             & LSTM & 0.9016 & 0.9272 \\ \hline
        \multirow{2}{*}{128} & CNN  & 0.8837 & 0.9101 \\ \cline{2-4} 
                             & LSTM & 0.9013 & 0.9286 \\ \hline
        \multirow{2}{*}{256} & CNN  & 0.8861 & 0.9091 \\ \cline{2-4} 
                             & LSTM & 0.9003 & 0.9313 \\ \hline
    \end{tabular}
    \caption{Model Performance Comparison across Maximum Sequence Lengths}
    \label{tab:seq_length_comparison}
\end{table}

For the CNN model, the highest accuracy for the dev set was at a maximum sequence length of 64, at 0.8888. The highest accuracy for the test set was at a maximum sequence length of 32, at 0.9133. So, the CNN model appears to perform best on the test set with shorter sequences, but needs more context to reach its peak on the dev set. For the LSTM model, the highest accuracy for the dev set was at a maximum sequence length of 16, at 0.9038. The highest accuracy for the test set was at a maximum sequence length of 256, at 0.9313. So, the LSTM model appears to perform best at a maximum length.

The results of this ablation study present a counterintuitive trend. Typically, an ablation study on sequence length is expected to highlight a significant degradation in model performance as the available context is reduced. However, our results demonstrate a surprising resilience and powerful performance at lower maximum sequence lengths. We identified the reasoning for this: In our preprocessing, we concatenate the title and description, which results in the most contextual information being processed regardless of sequence length. So, by our own preprocessing decisions, the effect of the ablation did not result in a significant decline in the performance of the two models. 

\section{Discussion}
This report aimed to investigate the CNN and LSTM model architectures' performance on the AG News dataset. Both models performed well and beat the baseline model's performance, as demonstrated by assignment one's results in Figure \ref{fig:assignment_1}. When examining the errors, most mistakes from the CNN can be traced back to specific words in different filter lengths of the sentence. It is expected that this can be helped by feeding the convolved signals into each other rather than pooling over all the filters. This way, the present context of different parts of the sentence can be combined more easily. This can be seen as the LSTM, which does contain multiple layers which transform the input to achieve better accuracy than the CNN. Despite these differences in accuracy the models do not differ wildly in the types of errors they make.

As for the ablation study, the results did not match the initial hypothesis. The data shows a strong resilience of the model to the shorter sequence lengths. We expect this is in part due to the way we structured the data preprocessing. As the data was preprocessed by concatenating the title and the description, while the sequences were cut by taking the first n tokens from the start, the shorter sequences ended up storing similar information. We expect that when changing the preprocessing to only include the description or when taking the sequence from the centre and not the start, that the hypothesised decrease in performance will be observed. 

\newpage
\printbibliography

\suppressfloats[t]
% Here be appendices!
\newpage
\appendix
\section{GitHub Repository}

The URL to the GitHub Repository: \href{https://github.com/TeunB6/natural-language-processing}{https://github.com/TeunB6/natural-language-processing}. For viewing purposes, the main branch is \texttt{main}.

\section{AI-Use Statement}

Generative AI was utilised with the purpose and intention of accelerating the development of the project with regard to providing examples on implementing certain features, for example, the CLI UX. Permitted tools were OpenAI ChatGPT, Google Gemini, and GitHub Copilot, with the agreement that training of the models was disabled. Generated code was only used as a reference for comparison to the original written code, to provide any insight on technical scope and considerations we may have missed during development. Autocomplete features were used with human review, and the final product of the project is 100\% human reviewed.

\section{Model Architectures}\label{ap:arch}
\begin{table}[h]
\centering
\caption{LSTM Classifier Model Architecture}
\label{tab:lstm-architecture}
\begin{tabular}{>{\bfseries}p{1cm} p{4cm} p{3cm} >{\raggedright\arraybackslash}p{5.5cm}}
\toprule
\textbf{Layer} & \textbf{Type} & \textbf{Output Shape} & \textbf{Parameters / Description} \\
\midrule
1 & Dropout & ($N$, $L$, $E$) & Dropout rate: 0.5 \\
\midrule
2 & LSTM & ($N$, $L$, $2H$) & 
\begin{tabular}[c]{@{}p{5.8cm}@{}}
- Input size: $E = 100$\\
- Hidden size: $H = 100$\\
- Number of layers: 4)\\
- Bidirectional: True\\
- Dropout between each layer with dropout rate: 0.5\\
- Output: final hidden states from all time steps
\end{tabular} \\
\midrule
3 & Hidden State Concatenation & ($N$, $2H$) & 
\begin{tabular}[c]{@{}p{5.8cm}@{}}
Concatenates last hidden states from both directions.\\
\end{tabular} \\
\midrule
4 & Dropout & ($N$, $2H$) & Dropout rate: 0.5 \\
\midrule
5 & Fully Connected (Linear) & ($N$, $C$) & 
\begin{tabular}[c]{@{}p{5.8cm}@{}}
- Input: $2H$\\
- Output: $C =4$ \\
- No activation (logits output)
\end{tabular} \\
\midrule
6 & Softmax (Optional) & ($N$, $C$) & 
\begin{tabular}[c]{@{}p{5.8cm}@{}}
Converts logits to class probabilities
\end{tabular} \\
\bottomrule
\end{tabular}
\vspace{0.5cm}
\footnotesize{\textbf{Note:} Here $N$ is the batch size and $L$ is the sequence length.}
\end{table}

\begin{table}[h]
\centering
\caption{CNN Classifier Model Architecture}
\label{tab:cnn-architecture}
\begin{tabular}{>{\bfseries}p{1cm} p{4cm} p{3cm} >{\raggedright\arraybackslash}p{5.5cm}}
\toprule
\textbf{Layer} & \textbf{Type} & \textbf{Output Shape} & \textbf{Parameters / Description} \\
\midrule
1 & Permute & ($N$, $E$, $L$) & 
\begin{tabular}[c]{@{}p{5.8cm}@{}}
Reshapes input from ($N$, $L$, $E$) to ($N$, $E$, $L$) for 1D Convolution
\end{tabular} \\
\midrule
2 & Convolutional Layers & ($N$, $F$, $L'$) & 
\begin{tabular}[c]{@{}p{5.8cm}@{}}
Multiple parallel convolutions with:\\
- Input channels: $E = 100$\\
- Output channels: $F = 100$\\
- Filter sizes: $K = [3,4,5,6,7]$\\
- Activation: ReLU after each convolution\\
- $L' = L - k + 1$ (sequence length after convolution, where $k\in K$)
\end{tabular} \\
\midrule
3 & Adaptive Max Pooling & ($N$, $F$, 1) & 
\begin{tabular}[c]{@{}p{5.8cm}@{}}
Global max pooling over the sequence dimension. Extracts the most important feature for each filter.
\end{tabular} \\
\midrule
4 & Squeeze & ($N$, $F$) & 
\begin{tabular}[c]{@{}p{5.8cm}@{}}
Removes the last dimension (size 1)
\end{tabular} \\
\midrule
5 & Concatenation & ($N$, $F \times |K|$) & 
\begin{tabular}[c]{@{}p{5.8cm}@{}}
Concatenates outputs from all filter sizes.\\
$F \times |K| = 100 \times 5 = 500$
\end{tabular} \\
\midrule
6 & Dropout & ($N$, $F \times |K|$) & Dropout rate: 0.5 \\
\midrule
7 & Fully Connected (Linear) & ($N$, $C$) & 
\begin{tabular}[c]{@{}p{5.8cm}@{}}
- Input: $F \times |K| = 500$\\
- Output: $C = 4$\\
\end{tabular} \\
\midrule
8 & Softmax (Optional) & ($N$, $C$) & 
\begin{tabular}[c]{@{}p{5.8cm}@{}}
Converts logits to class probabilities
\end{tabular} \\
\bottomrule
\end{tabular}

\vspace{0.5cm}
\footnotesize{\textbf{Note:} Here $N$ is the batch size and $L$ is the sequence length.}
\end{table}


\section{Cases of misclassification for the CNN}
\begin{table}[H]
\begin{longtable}{| m{2em} | m{25em} | m{5em} | m{5em} |}
        \hline
        \textbf{Ref} & \textbf{Text} & \textbf{Actual} & \textbf{Predicted} \\ \hline
        1 & Eleusis wins Long Island 'cap; Personal Legend wins Turnback the Eleusis made a successful US debut by beating Literacy by 2 lengths in Aqueduct's Saturday feature, the Grade II, \$150,000 Long Island Handicap for fillies and mares 3 and older. & Sports & Sci/Tech \\ \hline
         
        2 & Mr. Chang, Strike a Pose Halfway around the world, standing virtually in the middle of the Pacific Ocean, the incomparable Timmy Chang is just days away from throwing his first pass of the season. From my tattered sofa, I will be watching him. I want you to watch him, too. & Sports & Sci/Tech \\ \hline
        
        3 & Leeds students figure high on working curve Undergraduates in the city earn more than 90 a week on average, just behind Glasgow, Cambridge and Cardiff. Their hard-earned cash is likely to be spent on looking good and socialising. & Business & Sports \\ \hline
        
        4 & Apology, refund from cruise line In a move almost unheard of in its industry, Norwegian Cruise Line has apologized for service problems during the Pride of Aloha's first two months of sailing around Hawaii, and is refunding a portion of the service charge to everyone who has cruised on board. & Business & Sci/Tech \\ \hline
        
        5 & Wisdom was main course WALTHAM -- He is an 87-year-old man with a cane and a cigar, and the clout of a king. & Sports & World \\ \hline
        
        6 & Next Big Hit at the Modern: Its Reopening The buzz over the greatly enlarged museum is expected to turn into a cacophony on Saturday. & World & Sci/Tech \\ \hline
        
        7 & Tyncastle sale gets official go ahead Of the votes received by proxy and from shareholders in the room at a stormy extraordinary general meeting last night, 62.5 were in favour of the resolution. & Sports & Business \\ \hline
        
    8 & Robby Gordon Planning to Proceed With Caution Robby Gordon plans to join the oil spills, the tire chunks, the sharp pieces of debris and the other typical racetrack hazards on Sunday. & Sports & Business \\ \hline
    
    9 & Gators may be hungover before Cocktail Party Like the bourbon and beer, emotions tend to spill over in this traditional grudge match along the St. Johns River they call "The World's Largest Outdoor Cocktail Party." & Sports & World \\ \hline
    
    10 & Australia in grip of water crisis There is a warning that some of Australia's major cities could run out of drinking water. & Sci/Tech & World \\ \hline

\caption{The 10 lowest confidence errors from the CNN model}
\label{tab:misclassificationCNN}
\end{longtable}
\end{table}


\section{Cases of misclassification for the LSTM}
\begin{table}[H]
\begin{longtable}{| m{2em} | m{25em} | m{5em} | m{5em} |}
        \hline
        \textbf{Ref} & \textbf{Text} & \textbf{Actual} & \textbf{Predicted} \\ \hline
        1 & San Diego fiscally sound, mayor says SAN DIEGO - In the wake of another downgrading of San Diego's credit rating, Mayor Dick Murphy today reassured the public that the city is fiscally sound. & Business & Sports \\ \hline
        
        2 & Feed-Tube Law Is Struck Down in Florida Case The court said that Gov. Jeb Bush violated separation of powers when he signed a law to keep Theresa Schiavo alive. & World & Business \\ \hline
        
        3 & Vote 2004 - a guide to the primary Get ready for the primary with the Herald-Tribunes special news section profiling all the federal, state and local candidates in races in Tuesday's election. & World & Sci/Tech \\ \hline
        
        4 & Visit to a Potemkin village Gongzhong does not resemble any Tibetan village in Tibet. It is a village more from Epcot Center in Walt Disney World. & World & Sci/Tech \\ \hline
        
        5 & New \$50 bill designed to counter counterfeits WASHINGTON - The green is still there, but with touches of blue, red and yellow. A stylized image of the Stars and Stripes now waves in the background. & Business & Sports \\ \hline
        
        6 & Cape Town F1 bid And they plan to build a multi-million-dollar race track here for the event, an official said today. "The plans for the bid are far advanced." & Sports & Business \\ \hline
        
        7 & The Kid Stays in His Home, Dressed in PJ's, With a Live Mike Robert Evans, the fabled Hollywood producer and man about town, will be enjoying his latest gig, as a satellite radio talk-show host, from the comfort of his home. & Sci/Tech & Business \\ \hline
        
        8 & Slowing population 'lacks funds' Rich countries are giving only half the amount they promised to help to slow world population growth, the UN says. & Sci/Tech & World \\ \hline
        
        9 & Coast Guard Shuts Gulf of Mexico Ports HOUSTON (Reuters) - The U.S. Coast Guard shut five ports on Wednesday in the Gulf of Mexico coast states of Alabama, Florida and Mississippi as Hurricane Ivan churned nearer. & Business & World \\ \hline
        
        10 & Travel Column: Australia Through Aboriginal Eyes This week's TravelWatch column profiles Anangu Tours, an Aborigine-owned tour company in Australia's Red Center. & Sci/Tech & Business \\ \hline

\caption{The 10 lowest confidence errors from the LSTM model}
\label{tab:misclassificationLSTM}
\end{longtable}
\end{table}


\section{Additional Plots}
\label{ap:plots}

\begin{figure}[h]
    \centering
    
    % Row 1
    \begin{subfigure}[b]{0.48\textwidth}
        \centering
        \includegraphics[width=\textwidth]{reports/assignment2/img/cnn_training_history_maxlen_16.png}
        \caption{Max Sequence Length: 16}
        \label{fig:loss_16}
    \end{subfigure}
    \hfill
    \begin{subfigure}[b]{0.48\textwidth}
        \centering
        \includegraphics[width=\textwidth]{reports/assignment2/img/cnn_training_history_maxlen_32.png}
        \caption{Max Sequence Length: 32}
        \label{fig:loss_32}
    \end{subfigure}

    \vspace{20pt} % Increased vertical spacing for clarity

    % Row 2
    \begin{subfigure}[b]{0.48\textwidth}
        \centering
        \includegraphics[width=\textwidth]{reports/assignment2/img/cnn_training_history_maxlen_64.png}
        \caption{Max Sequence Length: 64}
        \label{fig:loss_64}
    \end{subfigure}
    \hfill
    \begin{subfigure}[b]{0.48\textwidth}
        \centering
        \includegraphics[width=\textwidth]{reports/assignment2/img/cnn_training_history_maxlen_128.png}
        \caption{Max Sequence Length: 128}
        \label{fig:loss_128}
    \end{subfigure}

    \vspace{20pt}

    % Row 3 - Centered single image
    \begin{subfigure}[b]{0.48\textwidth}
        \centering
        \includegraphics[width=\textwidth]{reports/assignment2/img/cnn_training_history_maxlen_256.png}
        \caption{Max Sequence Length: 256}
        \label{fig:loss_256}
    \end{subfigure}

    \caption{Training History for CNN models across all tested max sequence lengths.}
    \label{fig:cnn_maxlen}
\end{figure}

\begin{figure}[h]
    \centering
    
    % Row 1
    \begin{subfigure}[b]{0.48\textwidth}
        \centering
        \includegraphics[width=\textwidth]{reports/assignment2/img/lstm_training_history_maxlen_16.png}
        \caption{Max Sequence Length: 16}
        \label{fig:loss_16}
    \end{subfigure}
    \hfill
    \begin{subfigure}[b]{0.48\textwidth}
        \centering
        \includegraphics[width=\textwidth]{reports/assignment2/img/lstm_training_history_maxlen_32.png}
        \caption{Max Sequence Length: 32}
        \label{fig:loss_32}
    \end{subfigure}

    \vspace{20pt} % Increased vertical spacing for clarity

    % Row 2
    \begin{subfigure}[b]{0.48\textwidth}
        \centering
        \includegraphics[width=\textwidth]{reports/assignment2/img/lstm_training_history_maxlen_64.png}
        \caption{Max Sequence Length: 64}
        \label{fig:loss_64}
    \end{subfigure}
    \hfill
    \begin{subfigure}[b]{0.48\textwidth}
        \centering
        \includegraphics[width=\textwidth]{reports/assignment2/img/lstm_training_history_maxlen_128.png}
        \caption{Max Sequence Length: 128}
        \label{fig:loss_128}
    \end{subfigure}

    \vspace{20pt}

    \begin{subfigure}[b]{0.48\textwidth}
        \centering
        \includegraphics[width=\textwidth]{reports/assignment2/img/lstm_training_history_maxlen_256.png}
        \caption{Max Sequence Length: 256}
        \label{fig:loss_256}
    \end{subfigure}

    \caption{Training History for LSTM models across all tested max sequence lengths.}
    \label{fig:lstm_maxlen}
\end{figure}

\section{Assignment 1 Results}

\begin{figure}[htbp]
    \centering
    \setlength{\tabcolsep}{3pt}
    
    \begin{subtable}[b]{0.48\textwidth}
        \centering
        \begin{tabular}{|l|c|c|}
        \hline
        \textbf{Dataset} & \textbf{Accuracy} & \textbf{Weighted F1-score} \\ \hline
        Dev Set  & 0.8748 & 0.8742 \\ \hline
        Test Set & 0.7470 & 0.7455 \\ \hline
        \end{tabular}
        \caption{Logistic Regression Performance}
        \label{subtab:lr_perf}
    \end{subtable}
    \hfill
    \begin{subtable}[b]{0.48\textwidth}
        \centering
        \begin{tabular}{|l|c|c|}
        \hline
        \textbf{Dataset} & \textbf{Accuracy} & \textbf{Weighted F1-score} \\ \hline
        Dev Set  & 0.8776 & 0.8769 \\ \hline
        Test Set & 0.7661 & 0.7645 \\ \hline
        \end{tabular}
        \caption{SVC Performance}
        \label{subtab:svc_perf}
    \end{subtable}

    \vspace{0.6cm}

    % --- Row 2: LR Confusion Matrix ---
    \begin{subtable}[b]{1\textwidth}
        \centering
        \begin{tabular}{|l|cccc|cccc|}
        \hline
         & \multicolumn{4}{c|}{\textbf{Dev Set}} & \multicolumn{4}{c|}{\textbf{Test Set}} \\ \hline
        \textbf{Pred \textbackslash{} Act} & \textbf{World} & \textbf{Sport} & \textbf{Business} & \textbf{Sci/Tech} & \textbf{World} & \textbf{Sport} & \textbf{Business} & \textbf{Sci/Tech} \\ \hline
        \textbf{World}    & 2539 & 189 & 146 & 102 & 1520 & 152 & 101 & 127 \\
        \textbf{Sports}   & 52   & 2694 & 23  & 20  & 129  & 1588 & 67  & 116 \\
        \textbf{Business} & 169  & 22   & 2584 & 264 & 222  & 133  & 1235 & 310 \\
        \textbf{Sci/Tech} & 119  & 64   & 333  & 2680 & 196  & 122  & 248  & 1334 \\ \hline
        \end{tabular}
        \caption{Logistic Regression Confusion Matrix}
        \label{subtab:lr_cm}
    \end{subtable}

    \vspace{0.6cm}

    % --- Row 3: SVC Confusion Matrix ---
    \begin{subtable}[b]{1\textwidth}
        \centering
        \begin{tabular}{|l|cccc|cccc|}
        \hline
         & \multicolumn{4}{c|}{\textbf{Dev Set}} & \multicolumn{4}{c|}{\textbf{Test Set}} \\ \hline
        \textbf{Pred \textbackslash{} Act} & \textbf{World} & \textbf{Sport} & \textbf{Business} & \textbf{Sci/Tech} & \textbf{World} & \textbf{Sport} & \textbf{Business} & \textbf{Sci/Tech} \\ \hline
        \textbf{World}    & 2515 & 202 & 166 & 93  & 1542 & 143 & 88  & 127 \\
        \textbf{Sports}   & 36   & 2719 & 16  & 18  & 101  & 1632 & 52  & 115 \\
        \textbf{Business} & 146  & 23   & 2598 & 272 & 203  & 133  & 1251 & 313 \\
        \textbf{Sci/Tech} & 128  & 59   & 310  & 2699 & 192  & 125  & 186  & 1397 \\ \hline
        \end{tabular}
        \caption{SVC Confusion Matrix}
        \label{subtab:svc_cm}
    \end{subtable}

    \caption{Metric Results from Assignment 1}
    \label{fig:assignment_1}
\end{figure}

\newpage
\end{document}